\chapter{Calibration: Eras, Anchors, Nulls, \& Operating Points}
\label{sec:calibration}

Calibration turns the detector's fixed arithmetic into per-channel policy data
for a validated stable era.  It does not retune the pilot to whatever line happens to
be strongest: the synthesized target remains the nominal ATSC pilot.  Calibration
identifies the latest stable era of each channel, localizes the residual offset
around that target, characterizes the empirical null and correlation structure,
calibrates the science transfer, and constructs the residual-score histograms
used to select an operating point.  Every estimator is allowed to refuse when the
archive lacks the population it requires.  Calendar boundaries and older eras
remain essential provenance and evidence, but they are not arguments to the
operating-point selector.  Station and instrument changes each define
boundaries.  Their intersection forms a \emph{stable era}: a contiguous
interval with one accepted station state and one accepted instrument state.
``Station epoch'' in the archive discussion names only the station-state side
of that partition.  Each section below produces one per-channel object, and
each object appears once, for all 23 channels, in the table the section names.
Figure~\ref{fig:calibration:flow} draws the whole path from products to verdict.

\begin{figure}[t]
\centering
\resizebox{\linewidth}{!}{% Editable source for Figure \ref{fig:calibration:flow}: from products to a verdict.
\begin{tikzpicture}[font=\small, x=1cm, y=1cm,
  st/.style={draw, rounded corners=2pt, align=center, inner sep=3.5pt, minimum height=13mm, text width=#1, line width=0.75pt},
  ex/.style={st=#1, fill=ppMeasured!8, draw=ppMeasured!90!black},
  md/.style={st=#1, fill=ppModel!10, draw=ppModel!90!black},
  cd/.style={st=#1, fill=ppConditional!10, draw=ppConditional!90!black},
  arr/.style={-{Stealth[length=1.8mm]}, thick, ppPending!85!black},
  ln/.style={thick, ppPending!85!black}]
\node[db, minimum width=30mm, text width=28mm, minimum height=12mm] (prod) at (1.6,0) {\textbf{v5 products}: per frame $P_r$, $S_r[f]$, PSD, validity, UTC, ids};
\node[ex=30mm] (era) at (6.2,0) {\textbf{era classifier} (\S\ref{sec:calibration:eras})\\calibration-free; monthly medians; persistence};
\node[ex=30mm] (cur) at (10.8,0) {\textbf{current era}\\one station state, one instrument state};
\node[ex=24mm] (anc) at (2.2,-2.6) {\textbf{anchor} $f_a$\\on-minus-off, or fallback};
\node[ex=24mm] (nul) at (5.6,-2.6) {\textbf{null} $+$ source\\bulk median, left scale, CCDF};
\node[ex=24mm] (flr) at (9.0,-2.6) {\textbf{floor}\\off era, quiet subset, or refused};
\node[ex=24mm] (tau) at (12.4,-2.6) {$\boldsymbol{\tau_c}$\\measured, bound, or refused};
\node[md=34mm] (hist) at (6.0,-5.2) {\textbf{histogram family} $\{\mathcal H_\rho\}$\\counts, $r_{\rm sys}$ and $r_{\rm var}$ totals per candidate $\eta$};
\node[md=30mm] (sel) at (10.8,-5.2) {\textbf{RFIsher selector}\\cost under the bias constraint $\to (\rho^\star,\eta^\star)$};
\node[cd=30mm] (hold) at (10.8,-8.0) {\textbf{blocked holdout}\\frozen pair on untouched blocks, once};
\node[cd=30mm] (verd) at (15.4,-8.0) {\textbf{verdict}\\recovery, excision, or measurement-bound};
\draw[arr] (prod) -- (era); \draw[arr] (era) -- (cur);
% distribution line from the current era to the four objects
\draw[ln] (cur.south) -- ++(0,-0.45) coordinate (d1);
\draw[ln] (d1 -| anc.north) -- (d1 -| tau.north);
\foreach \t in {anc,nul,flr,tau} \draw[arr] (d1 -| \t.north) -- (\t.north);
% collector line from the four objects into the histogram family
\draw[ln] (anc.south) -- ++(0,-0.45) coordinate (c1);
\draw[ln] (c1 -| anc.south) -- (c1 -| tau.south);
\foreach \t in {nul,flr,tau} \draw[ln] (\t.south) -- (c1 -| \t.south);
\draw[arr] (c1 -| hist.north) -- (hist.north);
\draw[arr] (hist) -- (sel); \draw[arr] (sel) -- (hold); \draw[arr] (hold) -- (verd);
\node[md=30mm, minimum height=10mm] (tol) at (15.4,-5.2) {\textbf{science tolerance} $r_{\rm tol}$ and \textbf{transfer model} $(g_{\rm var}, g_{\rm sys})$};
\draw[arr] (tol) -- (sel);
\node[db, minimum width=28mm, text width=26mm, minimum height=9mm] (evalb) at (5.6,-8.0) {\textbf{evaluation blocks} (untouched)};
\draw[arr] (evalb) -- (hold);
\node[font=\scriptsize, text=ppFailure!90!black, align=left, text width=34mm, anchor=west] at (14.2,-2.6) {any refused input refuses the histogram; the selector never invents a floor or a $\tau_c$};
\end{tikzpicture}
}
\caption[From products to verdict]{From products to verdict. The era classifier is calibration-free and runs first; the four calibration objects are measured on the era it fixes; the histogram family is authorized only when none of them is refused; RFIsher selects one operating point; the blocked holdout evaluates it once on untouched blocks; and every arrow can terminate in a refusal that is reported rather than filled.}
\label{fig:calibration:flow}
\end{figure}

\section{Stable-Era Identification}
\label{sec:calibration:eras}

Every fraction the archive can quote averages over a span within which the
broadcast environment measurably changed --- the FCC incentive-auction repack
was still relocating stations through mid-2020, and individual transmitters
came and went throughout.  An archive-average operating point can therefore
describe no physical era: a station departure can move a channel from the
occupancy wall into a usable region, while a switch-on can move another channel
in the opposite direction, and a calibration that blends the two describes a
transmitter that never existed.  The unit of calibration is the channel's
\emph{current} era, chosen because it is the one that best represents the
channel's future.

The era is read from the fine-statistic time--frequency record.  For each
channel the survey products supply the fine statistic $T[f]$ of every valid
frame; binned by UTC month, that is a spectrogram of the detector statistic
across the fine span (Appendix~\ref{app:archive-diagnostics} shows one per
channel).  Thresholding the monthly median of $10\log_{10}(F/\mu_0)$ --- the
calibration-free coarse ratio --- marks every month as proxy-high or
proxy-low, and a change in the fine-peak location within the \emph{nominal}
pilot neighbourhood between proxy-high months marks a station change; the
procedure below is the authoritative statement of both.  A transition is placed at the boundary
between the last month in one state and the first in the next, or inside a
collection gap when the archive supplies no frames between them; the latest
interval with one accepted station state and one accepted instrument state is
the \emph{latest archive-observed stable era}, written ``current era'' for
short throughout.  The shorthand is a claim about the archive, not about the
present: a channel whose latest era ends well before the snapshot --- channel
30, whose collection ended in \rerun{September 2023} --- is marked
\emph{stale-latest} in Table~\ref{tab:calibration:eras}, and its era is called
current, or representative of the channel's future, only when resumed
monitoring or live data support that claim.  The rule is predeclared, calibration-free, and deliberately simple, so
that the era boundary is reproducible from the product and auditable by eye on
the plate, and so that it cannot be circular: the features are the monthly
fraction of valid frames raising the survey flag of
Eq.~\eqref{eq:survey-bootstrap-rule}, the monthly median pilot excess in the
nominal neighbourhood, the located anchor, the reference level, and the
recorded instrument state; a month enters only with at least thirty valid
frames; a state must persist for at least two consecutive populated months to
count as a transition rather than an excursion; a transition falling inside a
collection gap is placed at the gap with the gap's width as its boundary
uncertainty; and where a station or instrument record exists the transition is
checked against it and the agreement recorded.  None of these inputs is the
tolerance-selected $\eta^\star$ or anything derived from it --- the era must be
fixed before the threshold that is estimated on it.  A formal change-point
detector on the same monthly features \cite{killick2012} is the natural
refinement and is recorded as the sensitivity check, not the rule.  The frozen
procedure is:

\begin{enumerate}[leftmargin=2em,itemsep=1pt]
\item \emph{Inputs, all calibration-free.} Per valid frame: UTC; the exact
normalized coarse ratio $F/\mu_0$ ($\mu_0$ from weight norms alone); the fine
peak location within the \emph{nominal} acquisition neighbourhood
$f_{\rm pred} \pm 30$ bins (never the era-calibrated anchor); the reference
level $P_1 + P_2$; and the recorded receiver software tag and input map.
\item \emph{Aggregation.} UTC calendar month.  A month is populated only with at
least $30$ valid frames drawn from at least $5$ distinct acquisitions on at
least $3$ distinct days; policy values, with sensitivity values $(3, 5, 10)$
acquisitions recorded.
\item \emph{State.} A populated month is \emph{proxy-high} if its median
$10\log_{10}(F/\mu_0)$ is at least $1$~dB, \emph{proxy-low} if at most $0.5$~dB, and
\emph{ambiguous} between (sensitivity values $0.5$, $1$, $2$~dB); ``on'' and
``off'' are used only where a station record confirms the transmitter state; the
positive-excess rule itself fires on about half of quiet frames, which is why a
median excess and not a fired fraction is the statistic.  A \emph{station
change} is a shift of the monthly peak location by at least $3$ fine bins
($36$~Hz) against the running era.  An \emph{instrument change} is a change in
the recorded software tag or input map.
\item \emph{Persistence.} A new state or anchor must hold for at least $2$
consecutive populated months; a single-month excursion is flagged, not a
transition.  Ambiguous months inherit the state of their populated neighbours
if those agree and open a transition otherwise.
\item \emph{Placement and uncertainty.} A transition is placed between the last
month of the old state and the first of the new; when unpopulated months lie
between them it is placed at that gap and the gap's span is the boundary
uncertainty.
\item \emph{Fallback.} A channel with no proxy-low month has its current era bounded
only by station and instrument changes; its null then comes from the fine
bulk, declared as a mixture, or is refused (Section~\ref{sec:calibration:nulls}).
\item \emph{Naming.} A boundary is a \emph{transmitter sign-on} or
\emph{sign-off} only where a station record confirms it; otherwise it is a
\emph{spectral-state transition}.
\item \emph{Independence.} None of the inputs is the selected $\eta^\star$, the
era-calibrated anchor, or the era-calibrated null; the null of
Section~\ref{sec:calibration:nulls} is estimated \emph{on} the era after the
era is fixed, and the check that it does not feed back is that repeating the
procedure with the calibrated anchor substituted for the nominal neighbourhood
changes no boundary.
\end{enumerate}  Sign-offs and
switch-ons on the archive are frequent enough that the split matters on a
material fraction of the band: the superseded products resolved dated sign-offs
on channels \rerun{19, 20, 26, 27, and 32}, a switch-on on channel \rerun{35},
a two-epoch persistent carrier on channel \rerun{17}, and a collection-curtailed
archive on channel \rerun{30}.

\stubtab{v5 archive rerun}{One row per channel 14--36: current-era start and end
(UTC month) with boundary uncertainty; the transition evidence that defines the
start (sign-on, sign-off, station change, instrument change, or archive start)
and its agreement with any station or instrument record; a stale-latest flag
where the era ends before the snapshot; earlier eras listed for provenance;
current-era valid-frame count; and current-era calendar coverage (months with
frames / months spanned).}{Stable-era identification for every
channel: the current era used by every downstream calibration and its
transition evidence.\label{tab:calibration:eras}}

\section{Anchor Localization}
\label{sec:calibration:anchors}

For fine bin $f$, let $T_f^{\mathrm{on}}$ and $T_f^{\mathrm{quiet}}$ denote the
sets of recorded fine statistics in coarse-detected and coarse-quiet frames of
the current era.  The primary anchor estimator is
\begin{equation}
  \widehat f_a
  = \operatorname*{arg\,max}_{f}
  \left[
    \operatorname{median}(T_f^{\mathrm{on}})
    - \operatorname{median}(T_f^{\mathrm{quiet}})
  \right].
  \label{eq:calibration:anchor}
\end{equation}
Static structure common to both cohorts cancels, so the contrast isolates the
feature associated with the transmitter state rather than simply returning the
largest fixed instrumental bin.  A channel with no defensible quiet cohort uses
a plain-median fallback and is labelled as such; that fallback is adequate to
place a designated window around an always-present line but does not create a
transmitter-off null.

The current-era time-averaged spectrum around each pilot
(Figure~\ref{fig:census:psd}) shows the same anchor from the other side --- the
dominant lobe inside the fine span --- together with everything else the coarse
channel carries within $\pm 15$~kHz: companions, sidelobe forests on the
crowded-field channels, the instrumental channel-centre line where the pilot
happens to sit near it, and the two reference placements in clean territory.  A
lobe read from an averaged spectrum is not the same estimand as the
on-minus-off anchor of Eq.~\eqref{eq:calibration:anchor}, and the two are
reported side by side.  Their comparison against the fine span is a containment check on the
tap length, not a proof that $K = 128$ is optimal: Figure~\ref{fig:calibration:containment}
and Table~\ref{tab:calibration:anchors} place every measured anchor against
the $\pm 1.526$~kHz span of $K = 128$ and the $\pm 0.763$~kHz span of
$K = 256$, so the reader can see directly which channels the next radix-2 step
would clip (Section~\ref{sec:param:tradeoffs}).
On the superseded products the anchors lay \rerun{hundreds of hertz to
$-1.37$~kHz} from the synthesized tone, every one inside the $K = 128$ span,
and the largest --- channel 32's --- outside the $K = 256$ span; calendar
splits also revealed station changes rather than gradual drift (channel 32's
anchor moved by \rerun{11} fine bins and channel 35's by \rerun{29} between
eras, and channel 17's shifted between its two on-epochs as the later
transmitter configuration came to dominate).  The correct calibration object is
therefore $(c, e)$, channel and era, not channel alone; the runtime bundle
carries one anchor per era, and a wide union window is a temporary
compatibility measure, not a substitute for an era-specific calibration.

\stubfig{v5 archive rerun}{All 23 channels on one axis: the current-era anchor
offset from synthesis with its block-bootstrap error bar, over shaded bands for
the $K = 64$, $128$, and $256$ spans, so containment is read directly; channels
with a previous era show its anchor as an open marker.}{Anchor containment
against the fine spans, all channels.\label{fig:calibration:containment}}

\stubtab{v5 archive rerun}{One row per channel: current-era anchor
$\widehat f_a$ (fine bin and Hz from synthesis) with its estimator (on-minus-off
or fallback median) and block-bootstrap uncertainty; the dominant lobe of the
current-era averaged spectrum (Hz from synthesis, dB over median); the
per-frame peak-offset distribution (median, 90th and 99th percentiles, from
the retained fine spectra rather than the average); for $K = 64$, $128$, and
$256$, the fraction of frames and of pilot-associated energy $E_K(c)$ inside the span (against $E_{\min}$ of Eq.~\eqref{eq:param:kstar}, naming the channel that binds $K^\star$),
the straddle loss at the measured anchor, the margin to the span edge, and any
reference-region contamination (these come from the retained per-frame
$23.8$~Hz power spectra, which cover all three spans; they are a spectral
containment analysis, not a rerun of the detector at the other two tap
lengths, which the products cannot supply); the disposition (supported, supported with
sentinel, unsupported); and, for channels with more than one era, the anchor
of the previous era.}{Measured anchors against the fine span: the empirical
justification of $K = 128$.\label{tab:calibration:anchors}}

\section{A Floor Is Not the Smallest Observed Value}
\label{sec:calibration:floor}

A sensitivity floor is a property of a specified null population and estimator.
It is not the minimum value observed in a mixture of quiet and contaminated
frames.  For a channel with verified quiet data, calibration reports the
empirical null distribution of the selected statistic, its sample size, the
block structure used for uncertainty, the tail quantiles relevant to the false
alarm requirement, and the corresponding shelf-level mapping.  A robust core
width and an ordinary standard deviation are both useful because a thin
contaminated tail can leave the former nearly ideal while inflating the latter.

Channels that are effectively always on are censored in the wrong direction for
an empirical off-state floor.  For them, the calibration can still report:

\begin{enumerate}[leftmargin=2em]
  \item an analytical or injected-signal sensitivity under the declared noise model;
  \item the distribution of retained-frame proxy levels under a chosen policy; and
  \item an upper bound on a quiet-state floor if an independently verified off era exists elsewhere.
\end{enumerate}

It cannot call the least contaminated retained frame an off-state measurement.
The tempting substitute is the smallest shelf among frames that were detected,
and it is wrong in a way that is easy to miss and expensive to keep: the
inferred shelf tends to zero as the detection statistic approaches its null
centre from above, so the minimum detected shelf is not a physical floor but
the weakest positive excess that happened to occur in the sample, and it falls
without bound as the detection count grows.  On the superseded products it
evaluated to \rerun{$-76$~dB} on one always-on channel against measured floors
of \rerun{$-45$ to $-49$~dB} on the channels that have them --- thirty
decibels of sampling artefact, which would have inverted that channel's verdict.
Preparation therefore refuses to authorize a residual-score histogram without
a measured floor rather than substituting one, and an explicit override is the
only way to stand a value in from elsewhere.  A dated transmitter-off era
supplies the most defensible floor of all, because the off state is
established by the era rather than by the absence of a detection; where the
current era is an off era, its 90th-percentile shelf is the floor.

\section{The Null Distribution of the Current Era}
\label{sec:calibration:nulls}

The null is calibrated empirically because the independent-feed Gaussian
model of Chapter~\ref{sec:pipeline} is only a baseline.  Feed covariance,
receiver gain structure, spectral leakage, and intermittent non-Gaussian
features change the effective width and tail, and they change them per channel
and per era.  For each channel the histogram of the current era's frames --- the
normalized coarse statistic $F/\mu_0$ and the fine designated-window statistic
$Z(\rho)$ of Section~\ref{sec:detection:rule} --- is the object from which the
null is read.  It is in general a mixture: a null bulk of transmitter-off or
below-sensitivity frames and an upper tail of detected frames, and a histogram
of the era is not by itself the null.  Every channel therefore states the
\emph{source} of its null explicitly, from the following list in order of
strength: a verified transmitter-off era; an independently identified quiet
subset (frames quiet by a criterion other than the statistic being
calibrated); a control or reference-bin surrogate; a synthetic or model-only
null; or null unavailable, in which case the estimator refuses.  Where the
null is read from the bulk of a mixture, that is declared as such: its centre
is the median of the bulk and its scale is the left-side deviation about that
median, so that the detected tail cannot pull the calibration upward, and the
mixture assumption --- that the bulk is the $H_0$ population --- is stated
beside the number rather than left implicit.  The histogram is shown with its
tail complementary distribution (CCDF) so the reader can see where the bulk
ends and the tail begins.  The calibration
reports the bulk centre against the packed-weight prediction $\mu_0$, the raw
and robust-core widths against the i.i.d.\ widths of
Section~\ref{sec:estimator:statistics}, the tail fraction beyond a stated
number of core widths, and the shelf-level mapping of those quantiles.

Two checks accompany every null.  The width factor --- the ratio of the
measured to the idealized width --- is the number that says how far the
receiver departs from the i.i.d.\ model, and it is reported in both raw and
robust-core forms because a thin contaminated tail can inflate the former by an
order of magnitude while leaving the latter near ideal.  The exchangeability
check is what the rank rule requires: on quiet frames, a bulk bin tested
against the rank of the remaining bulk must exceed it at the combinatorial rate
$(|\mathcal B|+1-\rho)/(|\mathcal B|+1)$ implied by that rank.  On the
superseded products channel 29 met it to three significant figures at every
tested rank; the check is repeated per channel on the current era, and its
final form is performed on held-out quiet blocks, so that agreement supports
the rank calibration for that cohort rather than being assumed from the name of
the method.

\stubtab{v5 archive rerun}{One row per channel: current-era frame count and
null-bulk count; the null source (off era, quiet subset, reference surrogate,
model-only, or refused) and, for a bulk-read null, the declared mixture
assumption; coarse null centre of $F/\mu_0$ against $1$ (the packed-weight prediction) and its
left-side scale; coarse and fine width factors, raw and robust-core; tail
fraction beyond three core widths; fine-axis exchangeability rate at the
selected rank against the combinatorial prediction; kept-frame or off-era floor
in dB with its population; and the histogram plate reference in
Appendix~\ref{app:archive-diagnostics}.}{Current-era null calibration for
every channel: the object RFIsher's threshold selection consumes.\label{tab:calibration:nulls}}

\section{Correlation Time and Refusal Semantics}
\label{sec:calibration:correlation}

The residual chain needs a coherence scale, but the archive contains dense
samples only within a short capture and sparse samples between captures.  The
estimator therefore uses same-day or same-era structure functions with
bootstrap resampling over sidereal-day blocks and reports one of three outcomes:

\begin{equation}
  \widehat\tau_c \in
  \{\text{measured interval},\; \text{one-sided bound},\; \text{refused}\}.
\end{equation}

A refusal is propagated as a conservative sidereal-day cap in the screening
forecast and is visually and numerically distinguished from a measurement.  On
the superseded products a handful of channels carried measured values
(\rerun{34, 45, 62, and 165 minutes}), two carried an upper bound of
\rerun{five minutes} set by the shortest resolved lag, and the rest refused.
Because a verdict can change by orders of magnitude across the cadence gap
between the capture duration and the acquisition spacing, a contiguous-scan
campaign is a required measurement on the refused channels rather than an
optional precision improvement.  The outcome for every channel is carried in
Table~\ref{tab:tolerance:channels}.

\section{Stable-Era Operating Points}
\label{sec:calibration:operating}

Within the current era, calibration rejects invalid frames, validates the
anchor, designated set, bulk, and empirical null, applies the measured
correlation and coherence model, and maps the proxy through the declared
science transfer.  It also applies stability and coverage gates before
constructing the histogram.  A failed gate or missing transfer produces a
refusal, not an extra selector parameter.

The uncertainty gate keeps acquisition or sidereal-day blocks whole, resamples
the two era halves separately, and evaluates the worst movement over the
complete candidate surface.  Its seed, replicate count, coverage floor, minimum
block count, upper interval, and acceptable cost and residual margins are
explicit, recorded inputs.

The resulting input is the family $\{\mathcal H_\rho\}$ of
Eq.~\eqref{eq:detection:histogram}.  Its count bins determine the total and kept
frame counts, while its additive $r_{\rm sys}$ and optional $r_{\rm var}$ totals
are already in science units.  Each histogram also records $|\mathcal B|$, so the
selector can validate the one-based $\rho$ and report
$q_\rho=\rho/(|\mathcal B|+1)$.  The selector, RFIsher's threshold stage,
receives only this histogram family and the channel science tolerance.  It
evaluates the forecast cost
\begin{equation}
  \mathcal{C}(\pi)=\frac{1+r_{\rm var}(\pi)}{1-f(\pi)}
\end{equation}
subject to the separate coherent-bias constraint
$r_{\rm sys}(\pi)\le r_{{\rm sys},\rm tol}$.  Candidates retaining fewer than
thirty frames are omitted.  Among points within 2\% of the minimum cost, the
selector orders by lower $r_{\rm sys}$, less masking, lower $\rho$, and lower
$\eta$.  On the fine axis the selected policy is
$(f_a,\mathcal D,\mathcal B,\rho^\star,\eta^\star_{q16})$, although only $\rho$
and $\eta$ vary inside this final calculation; the rank is an optional second
degree of freedom, and a channel whose cost surface is flat in $\rho$ reports
the lowest rank in the plateau.  Until the visibility transfer is measured, the
survey reports $r_{\rm proxy}$ and carries $r_{\rm var}=g_{\rm var}r_{\rm proxy}$
and $r_{\rm sys}=g_{\rm sys}r_{\rm proxy}$ as explicit provisional closures.

For the final campaign, the calibration blocks build and select from the
accepted histogram, the chosen bundle is frozen, and a disjoint blocked holdout
is then evaluated without retuning.  This separates upstream stability evidence
from the threshold calculation and prevents held-out data from influencing the
selected pair.  A future runtime bundle carries an era validity interval and
refuses to load a calibration outside it; that interval is a deployment guard
and provenance field, not an input to the threshold calculation.  The selected
$(\rho^\star,\eta^\star)$ for every channel, with its cost, masked fraction, and
residual, is Table~\ref{tab:tolerance:eta}.

\section{Closing Calibration Campaign}
\label{sec:calibration:closing}

The final campaign produces one immutable evidence record per channel--stable-era
pair and marks exactly one, the latest accepted era, as operational.  Each record
contains anchor and uncertainty; designated and bulk masks; null sample count and
block definition; false-alarm curve; accepted residual-score histogram; selected
rank and Q16 multiplier where applicable; masked fraction and residual on disjoint
holdout blocks; correlation-time outcome; proxy-transfer version; era provenance;
and all source/product hashes.  Older records and their calendar boundaries remain
available for audit and state-change studies but are never blended into the current
selector input.  The campaign also includes simultaneous spectra
across each 6-MHz allocation and contaminated complex visibilities before and
after the sidereal filter.  Those latter products close the two transfer gates
that a pilot-only calibration cannot close by itself.
